\documentclass[10pt,a4paper]{amsart}
\usepackage[margin=19mm]{geometry}
\usepackage{amsmath,amssymb,mathtools}
\usepackage[scr=rsfs,cal=euler]{mathalfa}
\usepackage{xfrac}
\usepackage[unicode,hidelinks,bookmarks=false]{hyperref}
\newcommand{\ie}{i.e.\ }
\newcommand{\cf}{cf.\ }
\newcommand{\resp}{respectively\ }
\newcommand{\Subsection}{\subsection}
\providecommand{\nobreakdash}{\nobreak\hbox{-}}
\setlength{\emergencystretch}{2em}
\multlinegap0pt
\usepackage{amssymb}
\newtheorem{proposition}{Proposition}
\renewcommand{\Im}{\operatorname{Im}}
\renewcommand{\Re}{\operatorname{Re}}
\title{Two-component reduced functional}
\author{Senping Luo, Juncheng Wei}
\date{Source: arXiv:2605.07580v1 (CC BY 4.0)}
\begin{document}
\maketitle

\noindent\emph{Source and licence.} Two-component reduced functional. Version arXiv:2605.07580v1 (CC BY 4.0), distributed by arXiv under the Creative Commons Attribution 4.0 International licence, http://creativecommons.org/licenses/by/4.0/, as recorded in the arXiv licence metadata for this exact version and shown on the abstract page https://arxiv.org/abs/2605.07580v1. Full licence notice: \texttt{licenses/arxiv-2605.07580-CC-BY-4.0.txt}.\par\smallskip

\noindent\textit{Editorial scope.} Counted unit: the article's proposition Prop31 (source0.tex lines 1144--1174). The article's complete original proof of it is delivered with it (source0.tex lines 1177--1221, one \texttt{proof} environment of 45 lines). No mathematics is written, repaired or re-derived: the statement and the proof are the article's own lines, in the article's own notation, and no retained line is substituted or re-flowed.  No further source-local context is needed: every cross-reference of every retained line resolves inside the two delivered spans.  Nothing was omitted from any retained span, so the package carries no omission record.  No retained line contains a citation, so the wrapper carries no bibliography and no bibliography entry.  No retained line contains a figure environment, an image-inclusion command or drawing code, so no image is delivered and the package declares no asset dependency.  Editorial formatting only, in addition: an AMS \texttt{amsart} preamble that restates the article's own declarations and macros used by the retained lines, namely the environments \texttt{proposition} on separate counters, together with the standard packages those lines need.  The retained environments print in this document as Proposition 1 in order of appearance, whereas the article prints the same items with the numbering of its own full document.  The complete native source of the article as distributed in the arXiv e-print for this exact version is bundled unchanged as \texttt{sources/200090/source0.tex}.
\begin{proposition}[A minimum principle]\label{Prop31}  Assume that $\mathcal{W}$ is modular invariant, i.e.,
 \begin{equation}\aligned\label{M999}
\mathcal{W}(\frac{az+b}{cz+d})=\mathcal{W}(z),\;\;\hbox{for all}\;\;\left(
                                                                      \begin{array}{cc}
                                                                        a & b \\
                                                                        c & d \\
                                                                      \end{array}
                                                                    \right)\in \hbox{SL}_2(\mathbb{Z}),
\endaligned\end{equation}
and
 \begin{equation}\aligned\nonumber
\mathcal{W}(-\overline{z})=\mathcal{W}(z).
\endaligned\end{equation}

If \begin{equation}\aligned\label{Cabc1}
&\frac{\partial}{\partial y}\mathcal{W}(z)>0,\;\;z=(x,y)\in[0,\frac{1}{2}]\times[a,\infty)\;\;\hbox{for some}\;\;a>\frac{\sqrt3}{2},\\
&\frac{\partial}{\partial x}\mathcal{W}(z)<0,\;\;z=(x,y)\in[0,\frac{1}{2}]\times[b,\infty)\;\;\hbox{for some}\;\;b<\frac{\sqrt3}{2},
\endaligned\end{equation}
and
 \begin{equation}\aligned\label{fff}
\frac{a}{\frac{1}{4}+a^2}\geq b.
\endaligned\end{equation}

Then
\begin{equation}\aligned\nonumber
\min_{z\in\mathbb{H}}\mathcal{W}(z)=\min_{z\in\overline{\mathcal{D}_{\mathcal{G}}}}\mathcal{W}(z)\;\;\hbox{is attained at}\;\; e^{i\frac{\pi}{3}}(\hbox{hexagonal point}).
\endaligned\end{equation}
Here $\mathcal{D}_{\mathcal{G}}$ is the fundamental domain corresponding to modular group $\hbox{SL}_2(\mathbb{Z})$, explicitly, $\mathcal{D}_{\mathcal{G}}=\{
z\in\mathbb{H}: |z|>1,\; 0<x<\frac{1}{2}
\}.$
\end{proposition}

\begin{proof}
By the first part of \eqref{Cabc1}, we have
\begin{equation}\aligned\nonumber
\min_{z\in\overline{\mathcal{D}_{\mathcal{G}}}}\mathcal{W}(z)=\min_{z\in\overline{\mathcal{D}_{\mathcal{G}}}\cap\{y\leq a\}}\mathcal{W}(z).
\endaligned\end{equation}
We then assume $\min_{z\in\overline{\mathcal{D}_{\mathcal{G}}}\cap\{y\leq a\}}\mathcal{W}(z)$ is attained at $z_1:=(x_1,y_1)$.
Then $y_1\leq a$. Since $b<a$, by the second part of \eqref{Cabc1}, we have
\begin{equation}\aligned\label{xxx}
x_1=\frac{1}{2}.
\endaligned\end{equation}
Taking $\left(
                                                                      \begin{array}{cc}
                                                                        0 & 1 \\
                                                                        -1 & 1 \\
                                                                      \end{array}
                                                                    \right)\in \hbox{SL}_2(\mathbb{Z})$,
                                                                    we define
$z_2:=\left(
                                                                      \begin{array}{cc}
                                                                        0 & 1 \\
                                                                        -1 & 1 \\
                                                                      \end{array}
                                                                    \right)z_1$, it follows that
\begin{equation}\aligned\nonumber
z_2=\frac{1}{1-z_1}=\frac{\frac{1}{2}}{\frac{1}{4}+y_1^2}+i \frac{y_1}{\frac{1}{4}+y_1^2}
 \endaligned\end{equation}
and
 \begin{equation}\aligned\nonumber
\mathcal{W}(z_2)=\mathcal{W}(z_1).
 \endaligned\end{equation}
 This implies that $z_2$ still attains the minimum of $\min_{z\in\overline{\mathcal{D}_{\mathcal{G}}}}\mathcal{W}(z)$.
 Now we need an elementary inequality, namely,
  \begin{equation}\aligned\nonumber
\frac{u_2}{\frac{1}{4}+u_2^2}\geq\frac{u_1}{\frac{1}{4}+u_1^2}\;\;\hbox{if}\;\;\frac{1}{2}\leq u_2\leq u_1.
 \endaligned\end{equation}
 From this inequality, one has $\Im(z_2)\geq b$. In fact,
  \begin{equation}\aligned\nonumber
\Im(z_2)=\frac{y_1}{\frac{1}{4}+y_1^2}\geq\frac{a}{\frac{1}{4}+a^2}\;\;\hbox{if}\;\;\frac{1}{2}\leq y_1\leq a.
 \endaligned\end{equation}


By \eqref{fff}, we have $z_2\in\overline{\mathcal{D}_{\mathcal{G}}}\cap\{y\geq b\}$. Still by the second part of \eqref{Cabc1} and $z_2$ is the minimum point, there must be $\Re(z_2)=\frac{1}{2}$, i.e., $\frac{\frac{1}{2}}{\frac{1}{4}+y_1^2}=\frac{1}{2}$. It yields that $y_1=\frac{\sqrt3}{2}$. This and \eqref{xxx} yield the result. These complete the proof.


\end{proof}

\end{document}
